\documentclass[fontsize=15pt]{jlreq}
%プリアンブル
\usepackage{amsmath}
\usepackage{mathtools}
\pagestyle{empty}
\usepackage[margin=10mm]{geometry}
\usepackage{tikz}
\usetikzlibrary{arrows.meta, calc}
\usepackage{here}

\newcommand{\m}[1]{\raisebox{0.1em}{\textcircled{\scriptsize #1}}}



%本文
\begin{document}
\begin{center}
{\LARGE {振り返り}} \\
-生態\m{2}- %単元ごとに変えるか消す
\end{center}
\vspace{1em}
\begin{flushright}
名前：\underline{\hspace{5cm}}
\end{flushright}

% 問題部分
\begin{enumerate}
  \item 生産者，消費者，分解者をそれぞれ説明せよ．また，それぞれの
  例を挙げよ．
  \begin{itemize}
    \item 生産者：\\\\
    \hspace{1.05cm}例： 
    \item 消費者：\\\\
    \hspace{1.05cm}例： 
    \item 分解者：\\\\
    \hspace{1.05cm}例： 
  \end{itemize}

  \vspace{2em}
  
  \item 菌類と細菌類の違いを増え方の観点から説明せよ．\\
  (\hspace{17cm})

  \vspace{2em}

  \item 地球の二酸化炭素濃度が上昇した原因の1つとして，石油や
  石炭の大量消費が挙げられる．この石油や石炭などの燃料を総称して何というか．\\
  (\hspace{3cm})

  \vspace{2em}

  \item 大気中の二酸化炭素濃度が増加すると地球の平均気温が上がる．
  このことから，二酸化炭素にはどのようなはたらきがあると考えられるか．\\
  \small\small ヒント：太陽から地球へとやってきた熱の一部は本来宇宙へと逃げる．
  これによって地球の温度の異常な上昇を防いでいる．\normalsize\\
  (\hspace{17cm})
\end{enumerate}

\end{document}